%% Introduction section

\section{Introduction}
\label{sec:introduction}
% P1: Problem Context - The rise of DeFi and security challenges
Decentralized Finance (DeFi) has emerged as a major application of blockchain technology, enabling permissionless financial services through smart contracts.
By late 2025, the total value locked (TVL) in DeFi protocols had exceeded \$230 billion, reflecting the substantial financial value managed by DeFi systems~\cite{defiTVL2025}.
This rapid growth has attracted a large user base, with tens of millions of users interacting with DeFi protocols worldwide.
The increasing scale of DeFi has been accompanied by significant security challenges.
Industry analyses report that losses from cryptocurrency hacks and exploits exceeded \$3.4 billion in 2025, with DeFi protocols representing a major attack surface~\cite{cryptoLoss2025,chainalysis2025report}.
These attacks span both traditional smart-contract vulnerabilities, such as reentrancy as seen in the DAO incident, and advanced economic exploits that leverage flash loans to manipulate prices and protocol states within a single transaction~\cite{qin2020flashloan}.
The immutable nature of smart contracts exacerbates these risks.
Once deployed, vulnerabilities cannot be patched retroactively, making effective testing and continuous post-deployment monitoring critical for mitigating security threats in DeFi systems.
 

% P2: Limitations of existing approaches
Existing approaches to smart contract security can be broadly categorized into four paradigms, each with inherent limitations. \textit{Static analysis tools} such as Slither~\cite{slither} and Mythril~\cite{mythril2018} use pattern matching and symbolic execution to identify known vulnerability patterns. While effective at detecting classic bugs such as reentrancy, these tools suffer from path explosion in complex contracts and often struggle to capture runtime behaviors or cross-protocol interactions arising from DeFi composability~\cite{defi2023}. \textit{Formal verification frameworks} like Certora Prover~\cite{certoraProver} and KEVM~\cite{Hildenbrandt2018KEVM} provide mathematical guarantees but require substantial manual effort to specify invariants and often struggle to model the dynamic, multi-protocol nature of DeFi. \textit{Rule-based dynamic detection tools} such as DeFiRanger~\cite{defiranger2021}, DeFiTainter~\cite{defitainter2023}, and DeFort~\cite{defort2024} analyze transaction traces to detect price-manipulation attacks. However, these approaches rely on predefined mathematical models of specific pricing mechanisms (e.g., constant-product AMMs), making them less effective for custom protocols with non-standard price models; DeFiScope reports that such custom price models account for 44.2\% of 95 real-world price manipulation attacks collected over the past three years~\cite{zhong2025detectingvariousdefiprice}. \textit{LLM-based approaches} represent a recent paradigm shift. TrxGNNBERT~\cite{TrxGNNBERT} leverages a graph-transformer language model to learn trace-level representations for DeFi fraud detection, while GPTScan~\cite{GPTScan} combines GPT with static analysis for logic-vulnerability detection and iAudit~\cite{iaudit2025} employs fine-tuned models with multi-agent debate for source-code auditing. However, these methods either focus primarily on pre-deployment source-code analysis and cannot detect runtime attacks, or incur high end-to-end overhead for trace construction and model inference, limiting practicality for real-time monitoring. 
LLM hallucinations can lead to plausible but incorrect vulnerability reports.
Moreover, parametric knowledge quickly becomes outdated as attack patterns evolve.
%They also suffer from LLM hallucinations---producing plausible but incorrect vulnerability reports---and from knowledge staleness, since parametric knowledge cannot adapt to rapidly evolving attack patterns without costly retraining.

% P3: Key Insight and Challenges
Our key insight, distilled from years of practical experience in smart-contract security, is that DeFi attacks, despite their diverse manifestations, often exhibit recurring behavioral patterns at the semantic level. For example, a price manipulation attack typically follows a ``borrow, manipulate, and profit'' sequence: obtaining large capital via flash loans, executing trades that artificially inflate or deflate token prices, and extracting profit before repaying the loan, all within a single atomic transaction. This observation suggests that historical attack cases, annotated with behavioral semantics, can serve as a valuable knowledge source to guide the detection of new attacks. 
To empirically ground this insight, we analyze approximately 200 real-world DeFi attack incidents from DeFiHackLabs~\cite{defihacklabs2024}.
Despite their apparent diversity, these attacks consistently fall into eight major semantic categories, with Business Logic Flaws, Price Manipulation, and Reentrancy-based call attacks dominating the incident landscape.
%To empirically substantiate this insight, we analyze about 200 real-world attack incidents from DeFiHackLabs~\cite{defihacklabs2024} and find that these attacks cluster into eight major categories. The top three categories (Business Logic Flaws, Price Manipulation, and Reentrancy \& Call Attacks) account for the majority of all incidents. 
These results provide quantitative evidence that attack patterns are reusable across different protocols.

However, realizing this insight poses three technical challenges: \textbf{(C1) Semantic Gap:} Raw execution traces consist of low-level EVM operations and function calls that are difficult to interpret. \textit{How can we bridge the gap between low-level traces and high-level attack semantics?} \textbf{(C2) Retrieval Precision:} Na\"ive similarity search may retrieve irrelevant or misleading historical cases, introducing noise that can degrade detection accuracy. \textit{How can we ensure that retrieved cases are truly relevant to the query transaction?} \textbf{(C3) Reasoning Reliability:} Even with relevant context, LLMs may hallucinate or fail to correctly apply retrieved knowledge. \textit{How can we guide the LLM to make reliable attack-or-benign decisions?}

% P4: Our Approach
To address these challenges, we propose \approach, a Retrieval-Augmented Generation framework for the DeFi transaction attack detection. At a high level, \approach{} comprises two main components: a \textit{Historical Knowledge Base} and a \textit{RAG-based Decision Engine}. For \textbf{C1}, we construct an incrementally expandable knowledge base by collecting execution traces from historical attack incidents, extracting behavioral summaries using LLM-based semantic summarization, and storing them as vector embeddings alongside structured metadata. This transforms opaque execution traces into interpretable attack patterns. For \textbf{C2}, we design a two-stage retrieval mechanism: the first stage performs vector similarity search to retrieve top-$k$ candidate cases; the second stage applies a dynamic similarity threshold ($\tau = 0.85$) to filter out low-confidence matches, ensuring that only highly relevant cases are used as context. For \textbf{C3}, we employ an adaptive prompting strategy: when high-similarity historical cases exist, we construct an augmented prompt that instructs the LLM to reason by analogy with the retrieved attack patterns; otherwise, the system falls back to zero-shot analysis based solely on the transaction data. This design enables \approach{} to leverage historical knowledge when available while maintaining detection capability for novel attack patterns.

% P5: Results Summary
We evaluate \approach{} on two datasets: \dsystem{}, comprising 728 transactions (364 attacks across 110+ attack types and 364 benign transactions) curated from DeFiHackLabs, and \done{}, a specialized benchmark of 95 price-manipulation attacks~\cite{zhong2025detectingvariousdefiprice}. On \dsystem{}, \approach{} achieves an F1 score of 0.824 with a 15-second average latency, significantly outperforming TrxGNNBERT~\cite{TrxGNNBERT} (F1: 0.766, latency: $\sim$5 minutes). On \done{}, \approach{} attains 95.60\% recall, surpassing DeFiScope (80.22\%), DeFort (54.95\%), and DeFiRanger (52.75\%) by substantial margins. Ablation studies confirm that the RAG framework contributes an absolute F1 gain of 0.586 over a vanilla LLM, while the two-stage retrieval mechanism adds 0.030. Furthermore, \approach{} demonstrates cost efficiency: processing the entire \dsystem{} dataset costs only \$3.86 with GPT-4.1-mini or \$0.26 with Llama-3.1-8B-Instruct.

% P6: Contributions
This paper makes the following contributions:
\begin{itemize}
    \item We present \approach, the first RAG-based framework for the DeFi transaction attack detection that leverages historical attack knowledge to enhance LLM reasoning, addressing the limitations of both rule-based and pure LLM approaches.
    \item We design a two-stage retrieval mechanism with dynamic similarity thresholding that balances recall and precision, ensuring high-quality context for LLM-based reasoning while filtering out potentially misleading cases.
    \item We conduct comprehensive experiments on two datasets spanning 110+ attack types, demonstrating that \approach{} achieves state-of-the-art detection performance with significantly lower latency and cost than existing approaches.
    \item We provide an in-depth analysis of \approach{}'s capability boundaries, identifying scenarios where the system excels (e.g., reentrancy, access control) and where it faces challenges (e.g., cryptographic exploits, multi-step cross-contract attacks), offering insights for future research directions.
\end{itemize}

\noindent\textbf{Roadmap.} Section~\ref{sec:background} introduces the background. Section~\ref{sec:approach} details the design of \approach{}, including knowledge base construction, two-stage retrieval, and adaptive prompting. Section~\ref{sec:evaluation} presents our experimental evaluation and ablation studies, Section~\ref{sec:related_work} demonstrate the related work and Section~\ref{sec:conclusion} concludes with limitations and future directions.